\documentclass[a4paper, 10pt]{article}
% Per-subject config for this repo. See vendor/template/course-config.tex.example
% for what each macro is used for.

\newcommand{\CourseCode}{2110515}
\newcommand{\CourseName}{Intro to Quantum Computing}
\newcommand{\StudentID}{6733172621}
\newcommand{\StudentName}{Patthadon Phengpinij}
\newcommand{\DefaultCollaborators}{Claude (for\,\LaTeX\,styling and grammar checking)}

\usepackage{../vendor/template/CEDT-Assignment-style}

\begin{document}
\subject
\asgntitle{2}{}{Week 2}{\StudentID\ \StudentName}{\DefaultCollaborators}

% ================================================================================ %
%                                    Problem 01                                    %
% ================================================================================ %
\begin{problem}
\textbf{Normalization by Hand}
\par Normalize vector \( \ket{\vec{v_1}} = 2\ket{\uparrow} + 1\ket{\downarrow} \) by hand.
\end{problem}

\begin{solution}
	To normalize the vector \( \ket{\vec{v_1}} = 2\ket{\uparrow} + 1\ket{\downarrow} \), we need to divide it by its norm.

	The norm of \( \ket{\vec{v_1}} \) is given by:
	\[
		\abs{\ket{\vec{v_1}}} = \sqrt{|2|^2 + |1|^2} = \sqrt{4 + 1} = \sqrt{5}
	\]
	Therefore, the normalized vector is:
	\[
		\ket{\vec{v_1}_{\text{normalized}}} = \frac{2}{\sqrt{5}} \ket{\uparrow} + \frac{1}{\sqrt{5}} \ket{\downarrow}
		\approx \begin{pmatrix}
			0.8944 \\
			0.4472
		\end{pmatrix}
	\]
\end{solution}
% ================================================================================ %

% ================================================================================ %
%                                    Problem 02                                    %
% ================================================================================ %
\begin{problem}
\textbf{Normalization Using Python}
\par Check using Google Colab. You can use the scikit-learn library to do so:
\begin{codingbox}
import numpy as np

v1 = np.array([[2], [1]])
print(v1)

from sklearn.preprocessing import normalize

v2 = v1 / np.linalg.norm(v1)
print(v2)
\end{codingbox}
\end{problem}

\begin{solution}
	After running the provided Python code, the output is:
	\[
		\ket{\vec{v_1}_{\text{normalized}}} \approx \begin{pmatrix}
			0.8944 \\
			0.4472
		\end{pmatrix}
	\]
\end{solution}
% ================================================================================ %

% ================================================================================ %
%                                    Problem 03                                    %
% ================================================================================ %
\begin{problem}
\textbf{Normalization by Hand with Complex Coefficient}
\par Normalize vector \( \ket{\vec{v_1}} = 2\ket{\uparrow} + (1 + i)\ket{\downarrow} \) by hand.
\end{problem}

\begin{solution}
	The norm of \( \ket{\vec{v_1}} = 2\ket{\uparrow} + (1 + i)\ket{\downarrow} \) is calculated as follows:
	\[
		\abs{\ket{\vec{v_1}}} = \sqrt{|2|^2 + |1 + i|^2} = \sqrt{4 + (1^2 + 1^2)} = \sqrt{4 + 2} = \sqrt{6}
	\]
	
	Thus, the normalized vector is:
	\[
		\ket{\vec{v_1}_{\text{normalized}}} = \frac{2}{\sqrt{6}} \ket{\uparrow} + \frac{1 + i}{\sqrt{6}} \ket{\downarrow}
		\approx \begin{pmatrix}
			0.8165 \\
			0.4082 + 0.4082i
		\end{pmatrix}
	\]
\end{solution}
% ================================================================================ %

% ================================================================================ %
%                                    Problem 04                                    %
% ================================================================================ %
\begin{problem}
\textbf{Normalization Using Python with Complex Coefficient}
\par Verify using Google Colab. Remember you need to use \( 1 + 1j \) to represent \( 1 + i \).
\end{problem}

\begin{solution}
	Using the following Python code:
\begin{codingbox}
v1 = np.array([[2], [1+1j]])
v2 = np.round(v1 / np.linalg.norm(v1), 4)
print(f"Normalized vector:\n{v2}")
\end{codingbox}

	After running the provided Python code, the output is:
	\[
		\ket{\vec{v_1}_{\text{normalized}}} \approx \begin{pmatrix}
			0.8165 \\
			0.4082 + 0.4082i
		\end{pmatrix}
	\]
\end{solution}
% ================================================================================ %

% ================================================================================ %
%                                    Problem 05                                    %
% ================================================================================ %
\begin{problem}
\textbf{Vectors on a 2D Plane}
\par On a 2D plane, draw the vectors \( \ket{\vec{x^\prime}} = \frac{1}{\sqrt{2}}\paren{ \ket{\hat{x}} + \ket{\hat{y}} } \) and \( \ket{\vec{y^\prime}} = \frac{1}{\sqrt{2}}\paren{ \ket{\hat{x}} - \ket{\hat{y}} } \).
Write the column forms of \( \ket{\vec{x^\prime}} \) and \( \ket{\vec{y^\prime}} \).
Show that \( \ket{\vec{x^\prime}} \) and \( \ket{\vec{y^\prime}} \) are orthonormal,
i.e. \( \braket{\vec{x^\prime}|\vec{x^\prime}} = \braket{\vec{y^\prime}|\vec{y^\prime}} = 1 \)
and \( \braket{\vec{x^\prime}|\vec{y^\prime}} = \braket{\vec{y^\prime}|\vec{x^\prime}} = 0 \) and .
You have derived a new orthonormal basis for a 2D plane.
But this is not surprising, right?
As \( \ket{\vec{x^\prime}} \) and \( \ket{\vec{y^\prime}} \) can be obtained by just rotating
\end{problem}

\begin{solution}
	The column forms of \( \ket{\vec{x^\prime}} \) and \( \ket{\vec{y^\prime}} \), using \( \ket{\hat{x}} = \begin{pmatrix} 1 \\ 0 \end{pmatrix} \) and \( \ket{\hat{y}} = \begin{pmatrix} 0 \\ 1 \end{pmatrix} \), are
	\[
		\ket{\vec{x^\prime}} = \frac{1}{\sqrt{2}} \begin{pmatrix} 1 \\ 1 \end{pmatrix}, \qquad
		\ket{\vec{y^\prime}} = \frac{1}{\sqrt{2}} \begin{pmatrix} 1 \\ -1 \end{pmatrix}
	\]

	\begin{center}
		\begin{tikzpicture}[scale=2.2, >=Stealth]
			\draw[->] (-0.2, 0) -- (1.3, 0) node[right] {\( \hat{x} \)};
			\draw[->] (0, -0.2) -- (0, 1.3) node[above] {\( \hat{y} \)};

			\draw[->, thick, blue] (0, 0) -- ({1/sqrt(2)}, {1/sqrt(2)}) node[above right] {\( \ket{\vec{x^\prime}} \)};
			\draw[->, thick, red] (0, 0) -- ({1/sqrt(2)}, {-1/sqrt(2)}) node[below right] {\( \ket{\vec{y^\prime}} \)};

			\draw[dashed, gray] (0.55, 0) arc (0:45:0.55);
			\node at (0.68, 0.15) {\small \( 45^\circ \)};
		\end{tikzpicture}
	\end{center}

	\textbf{Norms:}
	\[
		\braket{\vec{x^\prime}|\vec{x^\prime}} = \frac{1}{2}\paren{1 \cdot 1 + 1 \cdot 1} = \frac{2}{2} = 1, \qquad
		\braket{\vec{y^\prime}|\vec{y^\prime}} = \frac{1}{2}\paren{1 \cdot 1 + (-1) \cdot (-1)} = \frac{2}{2} = 1
	\]

	\textbf{Cross terms:}
	\[
		\braket{\vec{x^\prime}|\vec{y^\prime}} = \frac{1}{2}\paren{1 \cdot 1 + 1 \cdot (-1)} = \frac{1 - 1}{2} = 0, \qquad
		\braket{\vec{y^\prime}|\vec{x^\prime}} = \frac{1}{2}\paren{1 \cdot 1 + (-1) \cdot 1} = \frac{1 - 1}{2} = 0
	\]

	Since both vectors have unit norm and their inner product with each other is zero,
	\( \ket{\vec{x^\prime}} \) and \( \ket{\vec{y^\prime}} \) form an orthonormal basis.

	This is not surprising: \( \ket{\vec{x^\prime}} \) and \( \ket{\vec{y^\prime}} \) are simply \( \ket{\hat{x}} \) and \( \ket{\hat{y}} \) rotated counterclockwise by \( 45^\circ \),
	and a rotation always preserves lengths and angles,
	so it necessarily maps an orthonormal basis to another orthonormal basis.
\end{solution}
% ================================================================================ %

% ================================================================================ %
%                                    Problem 06                                    %
% ================================================================================ %
\begin{problem}
\textbf{Basis Change by Hand}
\par Vector \( \vec{v_1} = \begin{pmatrix}
	1 \\ 2
\end{pmatrix} \) is in the \( \ket{+}/\ket{-} \) basis.
Represent it in the \( \ket{0}/\ket{1} \) basis.
Try to convert \( \ket{+} \) and \( \ket{-} \) in terms of \( \ket{0} \) and \( \ket{1} \) and perform the substitution.
To check your answer, convert it back to \( \ket{+}/\ket{-} \) using the equations introduced in this chapter and you should get back \( \begin{pmatrix}
	1 \\ 2
\end{pmatrix} \).
\end{problem}

\begin{solution}
	We have \( \ket{+} = \frac{1}{\sqrt{2}}\paren{\ket{0} + \ket{1}} \) and \( \ket{-} = \frac{1}{\sqrt{2}}\paren{\ket{0} - \ket{1}} \).
	Substituting \( \vec{v_1} = 1 \ket{+} + 2 \ket{-} \):
	\begin{align*}
		\vec{v_1} & =  1 \ket{+} + 2 \ket{-} \\
		 & = 1 \cdot \frac{1}{\sqrt{2}}\paren{\ket{0} + \ket{1}} + 2 \cdot \frac{1}{\sqrt{2}}\paren{\ket{0} - \ket{1}}    \\
		          & = \frac{1}{\sqrt{2}}\sqbracket{(1 + 2)\ket{0} + (1 - 2)\ket{1}}                                              \\
		\vec{v_1} & = \frac{3}{\sqrt{2}} \ket{0} - \frac{1}{\sqrt{2}} \ket{1}
	\end{align*}

	So in the \( \ket{0}/\ket{1} \) basis:
	\[
		\vec{v_1} = \begin{pmatrix}
			3/\sqrt{2} \\ -1/\sqrt{2}
		\end{pmatrix} \approx \begin{pmatrix}
			2.1212 \\ -0.7071
		\end{pmatrix}
	\]

	\textbf{Check.} Converting back to the \( \ket{+}/\ket{-} \) basis with \( \ket{0} = \frac{1}{\sqrt{2}}\paren{\ket{+} + \ket{-}} \) and \( \ket{1} = \frac{1}{\sqrt{2}}\paren{\ket{+} - \ket{-}} \), and writing \( \vec{v_1} = a\ket{0} + b\ket{1} \) with \( a = 3/\sqrt{2},\, b = -1/\sqrt{2} \):
	\[
		c_+ = \frac{a + b}{\sqrt{2}} = \frac{3/\sqrt{2} - 1/\sqrt{2}}{\sqrt{2}} = \frac{2/\sqrt{2}}{\sqrt{2}} = 1
	\]
	\[
		c_- = \frac{a - b}{\sqrt{2}} = \frac{3/\sqrt{2} + 1/\sqrt{2}}{\sqrt{2}} = \frac{4/\sqrt{2}}{\sqrt{2}} = 2
	\]
	which recovers \( \vec{v_1} = \begin{pmatrix} 1 \\ 2 \end{pmatrix} \) in the \( \ket{+}/\ket{-} \) basis, exactly as expected.
\end{solution}
% ================================================================================ %

% ================================================================================ %
%                                    Problem 07                                    %
% ================================================================================ %
\begin{problem}
\textbf{Vector Normalization by Hand and Using Python}
\par Normalize vector \( \ket{\vec{v_1}} = \begin{pmatrix}
	1 \\ 2
\end{pmatrix} \) (which is in the \( \ket{0}/\ket{1} \) basis) in the \( \ket{+}/\ket{-} \) basis.
Use Google Colab to verify your answer.
\end{problem}

\begin{solution}
	From Problem 6, \( \vec{v_1} \) expressed in the \( \ket{+}/\ket{-} \) basis (before normalization) is
	\[
		\vec{v_1} = \begin{pmatrix}
			3/\sqrt{2} \\ -1/\sqrt{2}
		\end{pmatrix}
	\]
	Its norm is
	\[
		\abs{\vec{v_1}} = \sqrt{\paren{\frac{3}{\sqrt{2}}}^2 + \paren{\frac{-1}{\sqrt{2}}}^2} = \sqrt{\frac{9}{2} + \frac{1}{2}} = \sqrt{5}
	\]
	so the normalized vector, in the \( \ket{+}/\ket{-} \) basis, is
	\[
		\ket{\vec{v_1}_{\text{normalized}}} = \frac{3/\sqrt{2}}{\sqrt{5}} \ket{+} - \frac{1/\sqrt{2}}{\sqrt{5}} \ket{-}
		= \frac{3}{\sqrt{10}} \ket{+} - \frac{1}{\sqrt{10}} \ket{-}
		\approx \begin{pmatrix}
			0.9487 \\ -0.3162
		\end{pmatrix}
	\]

	\newpage

	Using Python to verify:
\begin{codingbox}
import numpy as np

# v1 = (1, 2) in the |0>/|1> basis
v1_01 = np.array([[1], [2]])

# basis change matrix from |0>/|1> to |+>/|-> coefficients
M = (1 / np.sqrt(2)) * np.array([[1, 1], [1, -1]])

v1_pm = np.round(M @ v1_01, 4)
print("v1 in +/- basis (unnormalized):", v1_pm.ravel())

v1_pm_normalized = np.round(v1_pm / np.linalg.norm(v1_pm), 4)
print("v1 in +/- basis (normalized):", v1_pm_normalized.ravel())
\end{codingbox}
	which prints \( \begin{pmatrix} 0.9487 \\ -0.3162 \end{pmatrix} \), matching the hand-calculated result.
\end{solution}
% ================================================================================ %

% ================================================================================ %
%                                    Problem 08                                    %
% ================================================================================ %
\begin{problem}
\textbf{Measurement Probability}
\par For the \( \vec{v_1} \) in the last question, what is the probability to find \( \ket{+} \) in the \( \ket{+}/\ket{-} \) basis measurement?
\end{problem}

\begin{solution}
	The probability of measuring \( \ket{+} \) is the squared magnitude of its coefficient in the normalized vector:
	\[
		P(\ket{+}) = \abs{\frac{3}{\sqrt{10}}}^2 = \frac{9}{10} = 0.9
	\]
\end{solution}
% ================================================================================ %

% ================================================================================ %
%                                    Problem 09                                    %
% ================================================================================ %
\begin{problem}
\textbf{Vector Normalization by Hand and Using Python 2}
\par Normalize \( \ket{\vec{v_1}} \) in the \( \ket{0}/\ket{1} \) basis.
What is the probability to find \( \ket{0} \) in the \( \ket{0}/\ket{1} \) basis measurement?
Use Google Colab to verify your answer.
\end{problem}

\begin{solution}
	The norm of \( \vec{v_1} = \begin{pmatrix} 1 \\ 2 \end{pmatrix} \) in the \( \ket{0}/\ket{1} \) basis is
	\[
		\abs{\vec{v_1}} = \sqrt{1^2 + 2^2} = \sqrt{5}
	\]
	so the normalized vector is
	\[
		\ket{\vec{v_1}_{\text{normalized}}} = \frac{1}{\sqrt{5}} \ket{0} + \frac{2}{\sqrt{5}} \ket{1}
		\approx \begin{pmatrix}
			0.4472 \\ 0.8944
		\end{pmatrix}
	\]
	The probability to find \( \ket{0} \) is
	\[
		P(\ket{0}) = \abs{\frac{1}{\sqrt{5}}}^2 = \frac{1}{5} = 0.2
	\]

	\newpage

	Using Python to verify:
\begin{codingbox}
import numpy as np

v1 = np.array([[1], [2]])
v1_normalized = v1 / np.linalg.norm(v1)
print("v1 normalized in 0/1 basis:", v1_normalized.ravel())

prob_0 = np.abs(v1_normalized[0, 0]) ** 2
print("P(|0>):", prob_0)
\end{codingbox}
	which prints \( \begin{pmatrix} 0.4472 \\ 0.8944 \end{pmatrix} \) and \( P(\ket{0}) = 0.2000 \), matching the hand-calculated result.
\end{solution}
% ================================================================================ %

% ================================================================================ %
%                                    Problem 10                                    %
% ================================================================================ %
\begin{problem}
\textbf{Basis Change and Measurement}
\par If I get \( \ket{1} \) in the measurement in Problem 9, what is the probability of getting \( \ket{+} \) if I measure again in the \( \ket{+}/\ket{-} \) basis measurement?
(hint: after getting \( \ket{1} \), how to represent it in the \( \ket{+}/\ket{-} \) basis?)
\end{problem}

\begin{solution}
	After the measurement collapses the state to \( \ket{1} = \begin{pmatrix} 0 \\ 1 \end{pmatrix} \) in the \( \ket{0}/\ket{1} \) basis,
	we convert it to the \( \ket{+}/\ket{-} \) basis using \( a = 0,\, b = 1 \):
	\[
		c_+ = \frac{a + b}{\sqrt{2}} = \frac{0 + 1}{\sqrt{2}} = \frac{1}{\sqrt{2}}, \qquad
		c_- = \frac{a - b}{\sqrt{2}} = \frac{0 - 1}{\sqrt{2}} = -\frac{1}{\sqrt{2}}
	\]
	so \( \ket{1} = \frac{1}{\sqrt{2}} \ket{+} - \frac{1}{\sqrt{2}} \ket{-} \), which is already normalized.
	The probability of getting \( \ket{+} \) is
	\[
		P(\ket{+}) = \abs{\frac{1}{\sqrt{2}}}^2 = \frac{1}{2} = 0.5
	\]
\end{solution}
% ================================================================================ %

% ================================================================================ %
%                                     Appendix                                     %
% ================================================================================ %

\myappendix{code/assignment-02-code.pdf}{Assignment 02 | Code Solution}

% ================================================================================ %

\end{document}